\section{Introduction}
Boundary measurements for a harmonic equation record the response to an
imposed boundary value. In the anisotropic Calder\'on problem the unknown
is a Riemannian metric, and the response is its Dirichlet-to-Neumann map.
A diffeomorphism fixing the boundary leaves this map unchanged. The
uniqueness question is whether this coordinate freedom accounts for
every ambiguity in smooth metrics when the dimension is at least three.

This article gives an exposition of the affirmative argument in OpenAI's
September 24, 2026 preprint \citep{OpenAI2026}. That preprint treats
measurements supported in any nonempty boundary patch and, as a
consequence, full boundary measurements. We include the patch proof,
then deduce the formulation with one ordinary zero-frequency boundary
map in Theorem~\ref{main:full}. The manifold is compact and connected;
its boundary need not be connected, and the manifold need not be
orientable. All metrics are smooth and positive definite up to the
boundary. The two metrics in the full boundary statement have the same
tangential boundary metric, so their ordinary normal-derivative maps
use the same boundary density. No further measurement is used.

The organizing question is how boundary agreement can be continued into
the interior without analyticity. The Poisson embedding approach of
Lassas, Liimatainen, and Salo \citep{LLS20} provides useful context:
harmonic functions generated by sources encode points and local
coordinates. Here we follow OpenAI's Green-function construction.
Boundary symbol recovery first identifies every metric jet on the
observed patch. A common cap then gives a region on which the two
metrics and source kernels agree.

The main difficulty is to continue this agreement while a Green-function
pole approaches a shrinking sphere. Qualitative unique continuation
alone does not provide the needed uniform bound. OpenAI's argument
continues the Green kernels on a product of the two manifolds and
represents evaluation by a density on a moving sphere. An angular
estimate bounds this density independently of the sphere's shrinking
radius. Finite harmonic Hessian tests then recover the local metric,
and an open-and-closed argument continues the resulting isometry
throughout the manifold. We give the angular spectral mapping argument
in Subsection~\ref{prelim:spectral} before it is used in that estimate.

The coverage of this exposition is the complete smooth uniqueness
argument and its classical analytic inputs. The cap construction,
product continuation, shrinking-sphere estimate, and globalization
are due to OpenAI. The cited preprint is a source argument; this article
does not assert independent human certification of it. The proof uses
positive densities throughout, so local coordinate arguments do not
impose a global orientation. The final section returns to the full
ordinary Dirichlet-to-Neumann map and fixes every boundary point.
